\documentclass[12pt]{article}
\usepackage[paperwidth=612bp,paperheight=792bp,margin=0bp,noheadfoot]{geometry}
\usepackage{fontspec}
\usepackage{xeCJK}
\usepackage{tikz}
\pagestyle{empty}
\setlength{\parindent}{0pt}
\setlength{\parskip}{0pt}
\IfFontExistsTF{TeX Gyre Termes}{\setmainfont{TeX Gyre Termes}}{\setmainfont{Times New Roman}}
\IfFontExistsTF{Noto Serif CJK SC}{\setCJKmainfont{Noto Serif CJK SC}}{\setCJKmainfont{Songti SC}}
\newcommand{\QuestionTitle}{人机实验1 H-C1：萧相国世家}
\newcommand{\DrawQuestionTitle}{%
\node[anchor=base west,inner sep=0pt,outer sep=0pt] at ([xshift=72bp,yshift=-34bp]current page.north west) {{\fontsize{14bp}{14bp}\selectfont\bfseries \QuestionTitle}};%
\node[anchor=base west,inner sep=0pt,outer sep=0pt] at ([xshift=72bp,yshift=-62bp]current page.north west) {{\fontsize{12bp}{12bp}\selectfont 用时：\underline{\hspace{80bp}}}};%
}
\begin{document}
% Auto-generated for 萧何 (53.10.1 to 53.12.1)
\thispagestyle{empty}
\begin{tikzpicture}[remember picture,overlay]
\DrawQuestionTitle
\node[anchor=base west,inner sep=0pt,outer sep=0pt] at ([xshift=72bp,yshift=-84bp]current page.north west) {{\fontsize{12bp}{12bp}\selectfont 【53.10.1】 汉十一年，陈豨反叛，高祖亲自率军到了邯郸。}};
\node[anchor=base west,inner sep=0pt,outer sep=0pt] at ([xshift=72bp,yshift=-100bp]current page.north west) {{\fontsize{12bp}{12bp}\selectfont 【53.10.2】 平叛尚未结束，淮阴侯韩信又在关中谋反，吕后采用萧何的计策，杀了淮阴}};
\node[anchor=base west,inner sep=0pt,outer sep=0pt] at ([xshift=72bp,yshift=-116bp]current page.north west) {{\fontsize{12bp}{12bp}\selectfont 侯，此事记载在《淮阴侯列传》中。}};
\node[anchor=base west,inner sep=0pt,outer sep=0pt] at ([xshift=72bp,yshift=-132bp]current page.north west) {{\fontsize{12bp}{12bp}\selectfont 【53.10.3】 高祖已经听说淮阴侯被杀，派遣使者拜丞相萧何为相国，加封五千户，并令}};
\node[anchor=base west,inner sep=0pt,outer sep=0pt] at ([xshift=72bp,yshift=-148bp]current page.north west) {{\fontsize{12bp}{12bp}\selectfont 五百名士卒、一名都尉做相国的卫队。}};
\node[anchor=base west,inner sep=0pt,outer sep=0pt] at ([xshift=72bp,yshift=-164bp]current page.north west) {{\fontsize{12bp}{12bp}\selectfont 【53.10.4】 为此许多人都来祝贺，唯独召平表示哀悼。}};
\node[anchor=base west,inner sep=0pt,outer sep=0pt] at ([xshift=72bp,yshift=-180bp]current page.north west) {{\fontsize{12bp}{12bp}\selectfont 【53.10.5】 召平原是秦朝的东陵侯。}};
\node[anchor=base west,inner sep=0pt,outer sep=0pt] at ([xshift=72bp,yshift=-196bp]current page.north west) {{\fontsize{12bp}{12bp}\selectfont 【53.10.6】 秦朝灭亡后，他沦为平民，家中贫穷，在长安城东种瓜。}};
\node[anchor=base west,inner sep=0pt,outer sep=0pt] at ([xshift=72bp,yshift=-212bp]current page.north west) {{\fontsize{12bp}{12bp}\selectfont 【53.10.7】 他种的瓜味道甜美，所以社会上的人称它为“东陵瓜”，这是根据召平的封}};
\node[anchor=base west,inner sep=0pt,outer sep=0pt] at ([xshift=72bp,yshift=-228bp]current page.north west) {{\fontsize{12bp}{12bp}\selectfont 号来命名的。}};
\node[anchor=base west,inner sep=0pt,outer sep=0pt] at ([xshift=72bp,yshift=-244bp]current page.north west) {{\fontsize{12bp}{12bp}\selectfont 【53.10.8】 召平对相国萧何说：“祸患从此开始了。皇上风吹日晒地统军在外，而您留}};
\node[anchor=base west,inner sep=0pt,outer sep=0pt] at ([xshift=72bp,yshift=-260bp]current page.north west) {{\fontsize{12bp}{12bp}\selectfont 守朝中，未遭战事之险，反而增加您的封邑并设置卫队，这是因为目前淮阴侯刚刚在京城}};
\node[anchor=base west,inner sep=0pt,outer sep=0pt] at ([xshift=72bp,yshift=-276bp]current page.north west) {{\fontsize{12bp}{12bp}\selectfont 谋反，对您的内心有所怀疑。设置卫队保护您，并非以此宠信您，希望您辞让封赏不受，}};
\node[anchor=base west,inner sep=0pt,outer sep=0pt] at ([xshift=72bp,yshift=-292bp]current page.north west) {{\fontsize{12bp}{12bp}\selectfont 把家产、资财全都捐助军队，那么皇上心里就会高兴。”}};
\node[anchor=base west,inner sep=0pt,outer sep=0pt] at ([xshift=72bp,yshift=-308bp]current page.north west) {{\fontsize{12bp}{12bp}\selectfont 【53.10.9】 萧相国听从了他的计谋。}};
\node[anchor=base west,inner sep=0pt,outer sep=0pt] at ([xshift=72bp,yshift=-324bp]current page.north west) {{\fontsize{12bp}{12bp}\selectfont 【53.10.10】 高帝果然非常欢喜。}};
\node[anchor=base west,inner sep=0pt,outer sep=0pt] at ([xshift=72bp,yshift=-340bp]current page.north west) {{\fontsize{12bp}{12bp}\selectfont 【53.11.1】 汉十二年的秋天，黥布反叛，高祖亲自率军征讨他，多次派人来询问萧相国}};
\node[anchor=base west,inner sep=0pt,outer sep=0pt] at ([xshift=72bp,yshift=-356bp]current page.north west) {{\fontsize{12bp}{12bp}\selectfont 在做什么。}};
\node[anchor=base west,inner sep=0pt,outer sep=0pt] at ([xshift=72bp,yshift=-372bp]current page.north west) {{\fontsize{12bp}{12bp}\selectfont 【53.11.2】 萧相国因为皇上在军中，就在后方安抚勉励百姓，把自己的家财全都捐助军}};
\node[anchor=base west,inner sep=0pt,outer sep=0pt] at ([xshift=72bp,yshift=-388bp]current page.north west) {{\fontsize{12bp}{12bp}\selectfont 队，和讨伐陈豨时一样。}};
\node[anchor=base west,inner sep=0pt,outer sep=0pt] at ([xshift=72bp,yshift=-404bp]current page.north west) {{\fontsize{12bp}{12bp}\selectfont 【53.11.3】 有一个门客劝告萧相国说：“您灭族的日子不远了。您位居相国，功劳数第}};
\node[anchor=base west,inner sep=0pt,outer sep=0pt] at ([xshift=72bp,yshift=-420bp]current page.north west) {{\fontsize{12bp}{12bp}\selectfont 一，还能够再加功吗？您当初进入关中就深得民心，至今十多年了，民众都亲附您，您还}};
\node[anchor=base west,inner sep=0pt,outer sep=0pt] at ([xshift=72bp,yshift=-436bp]current page.north west) {{\fontsize{12bp}{12bp}\selectfont 是那么勤勉地做事，与百姓关系和谐，受到爱戴。皇上之所以屡次询问您的情况，是害怕}};
\node[anchor=base west,inner sep=0pt,outer sep=0pt] at ([xshift=72bp,yshift=-452bp]current page.north west) {{\fontsize{12bp}{12bp}\selectfont 您震撼关中。如今您何不多买田地，采取低价、赊借等手段来败坏自己的声誉？这样，皇}};
\node[anchor=base west,inner sep=0pt,outer sep=0pt] at ([xshift=72bp,yshift=-468bp]current page.north west) {{\fontsize{12bp}{12bp}\selectfont 上的心才会安定。”}};
\node[anchor=base west,inner sep=0pt,outer sep=0pt] at ([xshift=72bp,yshift=-484bp]current page.north west) {{\fontsize{12bp}{12bp}\selectfont 【53.11.4】 于是萧相国听从了他的计谋，高祖才非常高兴。}};
\node[anchor=base west,inner sep=0pt,outer sep=0pt] at ([xshift=72bp,yshift=-500bp]current page.north west) {{\fontsize{12bp}{12bp}\selectfont 【53.12.1】 高祖征罢黥布军队回来，民众拦路上书，说相国低价强买百姓田地房屋数量}};
\node[anchor=base west,inner sep=0pt,outer sep=0pt] at ([xshift=72bp,yshift=-516bp]current page.north west) {{\fontsize{12bp}{12bp}\selectfont 极多。}};
\end{tikzpicture}
\null
\end{document}

